\documentclass[
  man,
  longtable,
  a4paper,
  nolmodern,
  notxfonts,
  notimes,
  colorlinks=true,linkcolor=blue,citecolor=blue,urlcolor=blue]{apa7}

\usepackage{amsmath}
\usepackage{amssymb}




\RequirePackage{longtable}
\RequirePackage{threeparttablex}

\makeatletter
\renewcommand{\paragraph}{\@startsection{paragraph}{4}{\parindent}%
	{0\baselineskip \@plus 0.2ex \@minus 0.2ex}%
	{-.5em}%
	{\normalfont\normalsize\bfseries\typesectitle}}

\renewcommand{\subparagraph}[1]{\@startsection{subparagraph}{5}{0.5em}%
	{0\baselineskip \@plus 0.2ex \@minus 0.2ex}%
	{-\z@\relax}%
	{\normalfont\normalsize\bfseries\itshape\hspace{\parindent}{#1}\textit{\addperi}}{\relax}}
\makeatother




\usepackage{longtable, booktabs, multirow, multicol, colortbl, hhline, caption, array, float, xpatch}
\usepackage{subcaption}


\renewcommand\thesubfigure{\Alph{subfigure}}
\setcounter{topnumber}{2}
\setcounter{bottomnumber}{2}
\setcounter{totalnumber}{4}
\renewcommand{\topfraction}{0.85}
\renewcommand{\bottomfraction}{0.85}
\renewcommand{\textfraction}{0.15}
\renewcommand{\floatpagefraction}{0.7}

\usepackage{tcolorbox}
\tcbuselibrary{listings,theorems, breakable, skins}
\usepackage{fontawesome5}

\definecolor{quarto-callout-color}{HTML}{909090}
\definecolor{quarto-callout-note-color}{HTML}{0758E5}
\definecolor{quarto-callout-important-color}{HTML}{CC1914}
\definecolor{quarto-callout-warning-color}{HTML}{EB9113}
\definecolor{quarto-callout-tip-color}{HTML}{00A047}
\definecolor{quarto-callout-caution-color}{HTML}{FC5300}
\definecolor{quarto-callout-color-frame}{HTML}{ACACAC}
\definecolor{quarto-callout-note-color-frame}{HTML}{4582EC}
\definecolor{quarto-callout-important-color-frame}{HTML}{D9534F}
\definecolor{quarto-callout-warning-color-frame}{HTML}{F0AD4E}
\definecolor{quarto-callout-tip-color-frame}{HTML}{02B875}
\definecolor{quarto-callout-caution-color-frame}{HTML}{FD7E14}

%\newlength\Oldarrayrulewidth
%\newlength\Oldtabcolsep


\usepackage{hyperref}




\providecommand{\tightlist}{%
  \setlength{\itemsep}{0pt}\setlength{\parskip}{0pt}}
\usepackage{longtable,booktabs,array}
\usepackage{calc} % for calculating minipage widths
% Correct order of tables after \paragraph or \subparagraph
\usepackage{etoolbox}
\makeatletter
\patchcmd\longtable{\par}{\if@noskipsec\mbox{}\fi\par}{}{}
\makeatother
% Allow footnotes in longtable head/foot
\IfFileExists{footnotehyper.sty}{\usepackage{footnotehyper}}{\usepackage{footnote}}
\makesavenoteenv{longtable}

\usepackage{graphicx}
\makeatletter
\newsavebox\pandoc@box
\newcommand*\pandocbounded[1]{% scales image to fit in text height/width
  \sbox\pandoc@box{#1}%
  \Gscale@div\@tempa{\textheight}{\dimexpr\ht\pandoc@box+\dp\pandoc@box\relax}%
  \Gscale@div\@tempb{\linewidth}{\wd\pandoc@box}%
  \ifdim\@tempb\p@<\@tempa\p@\let\@tempa\@tempb\fi% select the smaller of both
  \ifdim\@tempa\p@<\p@\scalebox{\@tempa}{\usebox\pandoc@box}%
  \else\usebox{\pandoc@box}%
  \fi%
}
% Set default figure placement to htbp
\def\fps@figure{htbp}
\makeatother


% definitions for citeproc citations
\NewDocumentCommand\citeproctext{}{}
\NewDocumentCommand\citeproc{mm}{%
  \begingroup\def\citeproctext{#2}\cite{#1}\endgroup}
\makeatletter
 % allow citations to break across lines
 \let\@cite@ofmt\@firstofone
 % avoid brackets around text for \cite:
 \def\@biblabel#1{}
 \def\@cite#1#2{{#1\if@tempswa , #2\fi}}
\makeatother
\newlength{\cslhangindent}
\setlength{\cslhangindent}{1.5em}
\newlength{\csllabelwidth}
\setlength{\csllabelwidth}{3em}
\newenvironment{CSLReferences}[2] % #1 hanging-indent, #2 entry-spacing
 {\begin{list}{}{%
  \setlength{\itemindent}{0pt}
  \setlength{\leftmargin}{0pt}
  \setlength{\parsep}{0pt}
  % turn on hanging indent if param 1 is 1
  \ifodd #1
   \setlength{\leftmargin}{\cslhangindent}
   \setlength{\itemindent}{-1\cslhangindent}
  \fi
  % set entry spacing
  \setlength{\itemsep}{#2\baselineskip}}}
 {\end{list}}
\usepackage{calc}
\newcommand{\CSLBlock}[1]{\hfill\break\parbox[t]{\linewidth}{\strut\ignorespaces#1\strut}}
\newcommand{\CSLLeftMargin}[1]{\parbox[t]{\csllabelwidth}{\strut#1\strut}}
\newcommand{\CSLRightInline}[1]{\parbox[t]{\linewidth - \csllabelwidth}{\strut#1\strut}}
\newcommand{\CSLIndent}[1]{\hspace{\cslhangindent}#1}





\usepackage{newtx}

\defaultfontfeatures{Scale=MatchLowercase}
\defaultfontfeatures[\rmfamily]{Ligatures=TeX,Scale=1}





\title{Mental Health Evaluation through Text Analysis: umbrella project
documentation}


\shorttitle{Mental Health Evaluation through Text Analysis}


\usepackage{etoolbox}






\author{Gabriel Bonnin}



\affiliation{
{Ruhr University Bochum, Germany}}




\leftheader{Bonnin}






\authornote{ 

\par{       }
\par{Correspondence concerning this article should be addressed
to Gabriel
Bonnin, Email: \href{mailto:gabriel.bonnin@ruhr-uni-bochum.de}{gabriel.bonnin@ruhr-uni-bochum.de}}
}

\makeatletter
\let\endoldlt\endlongtable
\def\endlongtable{
\hline
\endoldlt
}
\makeatother
\RequirePackage{longtable}
\DeclareDelayedFloatFlavor{longtable}{table}

\urlstyle{same}



\makeatletter
\@ifpackageloaded{tcolorbox}{}{\usepackage[skins,breakable]{tcolorbox}}
\@ifpackageloaded{fontawesome5}{}{\usepackage{fontawesome5}}
\definecolor{quarto-callout-color}{HTML}{909090}
\definecolor{quarto-callout-note-color}{HTML}{0758E5}
\definecolor{quarto-callout-important-color}{HTML}{CC1914}
\definecolor{quarto-callout-warning-color}{HTML}{EB9113}
\definecolor{quarto-callout-tip-color}{HTML}{00A047}
\definecolor{quarto-callout-caution-color}{HTML}{FC5300}
\definecolor{quarto-callout-color-frame}{HTML}{acacac}
\definecolor{quarto-callout-note-color-frame}{HTML}{4582ec}
\definecolor{quarto-callout-important-color-frame}{HTML}{d9534f}
\definecolor{quarto-callout-warning-color-frame}{HTML}{f0ad4e}
\definecolor{quarto-callout-tip-color-frame}{HTML}{02b875}
\definecolor{quarto-callout-caution-color-frame}{HTML}{fd7e14}
\makeatother
\makeatletter
\@ifpackageloaded{caption}{}{\usepackage{caption}}
\AtBeginDocument{%
\ifdefined\contentsname
  \renewcommand*\contentsname{Table of contents}
\else
  \newcommand\contentsname{Table of contents}
\fi
\ifdefined\listfigurename
  \renewcommand*\listfigurename{List of Figures}
\else
  \newcommand\listfigurename{List of Figures}
\fi
\ifdefined\listtablename
  \renewcommand*\listtablename{List of Tables}
\else
  \newcommand\listtablename{List of Tables}
\fi
\ifdefined\figurename
  \renewcommand*\figurename{Figure}
\else
  \newcommand\figurename{Figure}
\fi
\ifdefined\tablename
  \renewcommand*\tablename{Table}
\else
  \newcommand\tablename{Table}
\fi
}
\@ifpackageloaded{float}{}{\usepackage{float}}
\floatstyle{ruled}
\@ifundefined{c@chapter}{\newfloat{codelisting}{h}{lop}}{\newfloat{codelisting}{h}{lop}[chapter]}
\floatname{codelisting}{Listing}
\newcommand*\listoflistings{\listof{codelisting}{List of Listings}}
\makeatother
\makeatletter
\makeatother
\makeatletter
\@ifpackageloaded{caption}{}{\usepackage{caption}}
\@ifpackageloaded{subcaption}{}{\usepackage{subcaption}}
\makeatother

% From https://tex.stackexchange.com/a/645996/211326
%%% apa7 doesn't want to add appendix section titles in the toc
%%% let's make it do it
\makeatletter
\xpatchcmd{\appendix}
  {\par}
  {\addcontentsline{toc}{section}{\@currentlabelname}\par}
  {}{}
\makeatother

%% Disable longtable counter
%% https://tex.stackexchange.com/a/248395/211326

\usepackage{etoolbox}

\makeatletter
\patchcmd{\LT@caption}
  {\bgroup}
  {\bgroup\global\LTpatch@captiontrue}
  {}{}
\patchcmd{\longtable}
  {\par}
  {\par\global\LTpatch@captionfalse}
  {}{}
\apptocmd{\endlongtable}
  {\ifLTpatch@caption\else\addtocounter{table}{-1}\fi}
  {}{}
\newif\ifLTpatch@caption
\makeatother

\begin{document}

\maketitle




\setlength\LTleft{0pt}




\begin{tcolorbox}[enhanced jigsaw, breakable, coltitle=black, bottomrule=.15mm, rightrule=.15mm, left=2mm, arc=.35mm, bottomtitle=1mm, colframe=quarto-callout-note-color-frame, toprule=.15mm, leftrule=.75mm, colbacktitle=quarto-callout-note-color!10!white, titlerule=0mm, colback=white, toptitle=1mm, title=\textcolor{quarto-callout-note-color}{\faInfo}\hspace{0.5em}{Note}, opacitybacktitle=0.6, opacityback=0]

\includegraphics[width=1.875in,height=\textheight,keepaspectratio]{meta_adult_manuscript_files/mediabag/reports/Logo_RUBRS.pdf}

This project is supported by the
\href{https://www.research-school.rub.de}{Ruhr University Bochum
Research School}.

\end{tcolorbox}

\section{Abstract}\label{abstract}

Psychotherapy is one of the most effective treatments for mental health
problems, but its success depends on accurate diagnostic assessments.
Current assessment practices largely rely on standardized closed-ended
scales that, while reliable, may fail to capture the complexity, context
and individuality of patients' mental states. Advances in artificial
intelligence (AI) and natural language processing (NLP) enable the
measurement of psychological constructs through natural language,
offering a promising complement to traditional assessment methods by
leveraging patients' own descriptions of their experiences.

While previous research has primarily focused on social media or
non-clinical language, this project applies state-of-the-art large
language models (LLMs) to open-ended intake data from a German
outpatient psychotherapy clinic, providing an ecologically valid
resource for studying language-based assessment in real-world clinical
settings.

The project comprises several complementary research areas. First, we
examine whether patient language at intake reflects cross-sectional
symptom severity and impairment and whether it provides incremental
information beyond established self-report questionnaires. Second, we
conduct a thematic analysis of patient responses to different open-ended
prompts to characterize recurring themes. Third, we evaluate whether
pre-therapy language predicts longitudinal treatment outcomes beyond
baseline symptom measures.

Together, these studies aim to clarify how patient-generated language
can be used for assessment, interpretation, and prognosis in
psychotherapy. The findings are intended to inform the development of
clinically meaningful, language-based assessment tools that support
personalized care and contribute to the modernization of mental health
evaluation.

\section{Introduction}\label{introduction}

Mental health problems pose a significant global challenge, accounting
for a considerable proportion of deaths and disability-adjusted life
years (\citeproc{ref-WorldHealthOrganization2017}{World Health
Organization, 2017}). Psychotherapy is an effective and sustainable
intervention for reducing symptoms and improving quality of life
(\citeproc{ref-chorpita2011evidence}{Chorpita et al., 2011};
\citeproc{ref-wampold2015great}{Wampold \& Imel, 2015}), but its success
critically depends on accurate assessments
(\citeproc{ref-Jensen-Doss2008}{Jensen-Doss \& Weisz, 2008};
\citeproc{ref-Lutz2022}{Lutz et al., 2022}).

Standardized closed-ended tools such as the Beck Depression Inventory-II
(\citeproc{ref-Beck1996}{Beck et al., 1996}) rely on numerical scales
(\citeproc{ref-Likert1932}{Likert, 1932}) to structure and standardize
assessments and are widely used in clinical research and practice. While
these methods have advanced replicability and reliability in
psychological assessment, they can miss important individual differences
by restricting responses to pre-defined categories, limiting the ability
to capture the complexity of mental states
(\citeproc{ref-Kjell2024}{Kjell, Kjell, et al., 2024}).

To address this limitation, researchers have turned to patients' natural
language, which provides high-dimensional, context-rich information that
remains underutilized in current assessment practices. Advances in AI,
particularly transformer-based LLMs (\citeproc{ref-Vaswani2017}{Vaswani
et al., 2017}), enable the context-rich analysis of natural language
with remarkable accuracy (\citeproc{ref-devlin-etal-2019-bert}{Devlin et
al., 2019}). When applying LLMs to mental health assessment, two primary
measurement contexts emerge. On the one hand, ecologically collected
speech data can be analysed that was not primarily generated for
measurement purposes. Examples include the prediction of depression
levels on the basis of social-media-posts
(\citeproc{ref-eichstaedt2018}{Eichstaedt et al., 2018}), or
psychotherapy session transcripts for analyzing emotional reaction
(\citeproc{ref-eberhardt2024}{Eberhardt et al., 2024};
\citeproc{ref-tanana2021}{Tanana et al., 2021}). On the other hand,
language-based self-reporting can be used specifically for measurement
purposes, for example in the form of standardized questions
(``\emph{Please describe how you have been feeling over the past two
weeks.}''). Empirical studies highlight the potential of NLP-based
analysis of prompted open-ended responses, achieving moderate
convergence with closed-ended rating scales using traditional NLP
methods (\citeproc{ref-kjell2019semantic}{Kjell et al., 2019}) and
nearing theoretical upper limits of accuracy with LLMs
(\citeproc{ref-Kjell2022}{Kjell et al., 2022}).

However, much of the existing literature relies on online survey data or
non-clinical text data, leaving a gap in our understanding of how these
methods perform in real-world clinical contexts. In contrast, open-ended
patient responses are routinely collected in clinical settings as part
of pre-therapy intake procedures but remain largely underused in
empirical research. The present project leverages data from an
outpatient psychotherapy clinic at Ruhr University Bochum, where
patients respond to multiple prompts addressing the development and
context of their mental health problems, perceived causes, social
reactions, current difficulties, and therapy goals. These narratives are
linked to structured diagnostic interviews, repeated symptom
assessments, clinician ratings, and longitudinal outcome measures. This
longitudinal clinical dataset enables a comprehensive examination of
patient language across multiple analytic perspectives.

Accordingly, the present project is organized into three complementary
research areas. The first investigates whether pre-therapy language
reflects cross-sectional symptom severity and clinician-rated impairment
and whether language-based representations provide incremental
information beyond standardized self-report questionnaires. The second
focuses on the semantic content and structure of pre-therapy patient
narratives, using question-specific analyses to identify recurring
themes and selective response patterns in how patients conceptualize
their mental health problems before therapy. The third evaluates the
prognostic value of pre-therapy language by testing whether patient
narratives predict treatment response and individualized goal attainment
over time, beyond baseline symptom severity.

By integrating assessment, interpretive, and prognostic perspectives,
this project aims to advance the clinical use of natural language in
psychotherapy. Ultimately, the findings seek to support more nuanced,
patient-centered assessment practices and contribute to the development
of language-based tools that complement existing diagnostic frameworks
and inform personalized treatment planning.

\section{Shared Methods}\label{shared-methods}

All substudies draw on the same clinical cohort and share a common set
of intake and outcome variables. The shared dataset comprises (a)
pre-therapy intake data, including sociodemographics, standardized
psychometric questionnaires, and question-specific open-ended patient
narratives, and (b) longitudinal clinical measures collected repeatedly
during and after treatment. Textual analyses are based exclusively on
pre-therapy narratives, while psychometric and clinician-rated measures
are used as cross-sectional outcomes, covariates, or longitudinal
endpoints depending on the substudy.

The following sections describe the shared dataset, preprocessing
pipeline, and measurement instruments used across all substudies.

\subsection{Measures}\label{measures}

\subsubsection{Sociodemographic and context
measures}\label{sociodemographic-and-context-measures}

Sociodemographic information included age, sex, marital and relationship
status, general educational attainment, vocational qualification, and
current work ability. Contextual variables captured prior psychological
or psychiatric treatment and the manner in which therapy ended (e.g.,
regular completion, dropout).

\subsubsection{Responses from open-ended questions before
therapy}\label{responses-from-open-ended-questions-before-therapy}

At the start of therapy, patients complete two separate questionnaires
designed to assess key aspects of their mental health concerns,
functional impairments, and expectations for treatment. Questions 1--9
come from the first questionnaire (\emph{Fragebogen zur
Lebensgeschichte}), and questions 10--12 come from the second
(\emph{Eingangsfragebogen}). The questions include:

\begin{enumerate}
\def\labelenumi{\arabic{enumi}.}
\tightlist
\item
  \textbf{Problem development:} `Briefly describe how the problems for
  which you are seeking treatment have developed over time.' (german
  original question: „Beschreiben Sie kurz, wie sich Ihre Probleme,
  wegen derer Sie eine Behandlung aufsuchen, im Laufe der Zeit
  entwickelt haben.'')
\item
  \textbf{Extra stressors:} `What causes you stress in addition to your
  everyday problems (e.g.~finances, housing situation)?' (german
  original question: „Was macht Ihnen zusätzlich zu Ihren Problemen im
  Alltag Stress (z. B. Finanzen, Wohnsituation)?{}``)
\item
  \textbf{Pre-onset changes:} `Did something special change in your life
  before the onset of your symptoms? (e.g.~death of an important person,
  divorce or separation, change in work situation or income, addition to
  the family)' (german original question: „Hat sich vor dem Beginn Ihrer
  Beschwerden etwas Besonderes in Ihrem Leben verändert? (z. B. Tod
  einer wichtigen Bezugsperson, Scheidung oder Trennung, Veränderung der
  Arbeitssituation oder des Einkommens, Familienzuwachs)``)
\item
  \textbf{Event connection:} `Do you see a connection between the
  event(s) and the development of your problems?' (german original
  question: „Sehen Sie einen Zusammenhang zwischen dem Ereignis/den
  Ereignissen und der Entwicklung Ihrer Probleme?{}``)
\item
  \textbf{Physical symptoms:} `Are there any physical side effects when
  your problems occur?' (german original question: „Gibt es körperliche
  Begleiterscheinungen, wenn Ihre Probleme auftreten?{}``)
\item
  \textbf{Problem causes:} `What do you think are the causes of your
  problems?' (german original question: „Welche Ursachen sehen Sie für
  Ihre Probleme?{}``)
\item
  \textbf{Expected improvements:} `What would improve in your life if
  you no longer had your problems?' (german original question: „Was
  würde sich in Ihrem Leben verbessern, wenn Sie ihre Probleme nicht
  mehr hätten?{}``)
\item
  \textbf{Environment response:} `How does your environment (partner,
  family, friends, work colleagues) react to your problems?' (german
  original question: „Wie reagiert Ihre Umwelt (Partner:in, Familie,
  Freund:innen, Arbeitskolleg:innen) auf die Probleme?{}``)
\item
  \textbf{No change required:} `What should not change under any
  circumstances as a result of the therapy?' (german original question:
  „Was sollte sich durch die Therapie auf keinen Fall verändern?{}``)
\item
  \textbf{Problem description:} `Finally, please describe in your own
  words the problems for which you would like treatment.' (german
  original question: „Beschreiben Sie zum Abschluss bitte noch einmal in
  eigenen Worten Ihre Probleme, deretwegen Sie eine Behandlung
  wünschen.``)
\item
  \textbf{Impacted life areas:} `In which areas of your life do these
  problems limit you (e.g.~job, relationship)?' (german original
  question: „In welchen Lebensbereichen schränken Sie diese Probleme ein
  (z. B. Beruf, Partnerschaft)?{}``)
\item
  \textbf{Therapy goals:} `What would you like to achieve for yourself
  in therapy?' (german original question: „Was möchten Sie in der
  Therapie für sich erreichen?{}``)
\end{enumerate}

\subsubsection{Psychometric measures}\label{psychometric-measures}

Clinical and psychometric variables were retrieved from the FBZ database
and included diagnostic information, self-report symptom measures,
therapist- and patient-rated outcome measures, positive mental health
indicators, and therapeutic process variables. Diagnoses were coded
according to DSM-5 and ICD-10 criteria. Symptom severity and treatment
outcomes were assessed using a combination of standardized self-report
questionnaires and clinician-rated instruments administered at different
points during treatment.

\paragraph{Diagnosis.}\label{diagnosis}

Diagnosis at the outpatient clinic is conducted using structured
clinical interviews. These typically take place before therapy begins,
usually at the fourth therapist--patient contact. The interview used is
the Diagnostic Interview for Mental Disorders
(\citeproc{ref-margraf2021}{Margraf et al., 2021}), which covers the
most frequent DSM-5 disorders encountered in outpatient therapy
settings.

\paragraph{Beck-Depression-Inventory
II.}\label{beck-depression-inventory-ii}

Depressive symptoms were assessed using the \emph{Beck Depression
Inventory--II} (BDI-II; (\citeproc{ref-Beck1996}{Beck et al., 1996})), a
widely used self-report questionnaire measuring the severity of
depressive symptomatology over the past two weeks.

\paragraph{Depression Anxiety Stress Scale
42.}\label{depression-anxiety-stress-scale-42}

Depressive symptoms, anxiety symptoms, and general psychological
distress were assessed using the \emph{Depression Anxiety Stress
Scale--42} (DASS-42(\citeproc{ref-Lovibond1995}{Lovibond \& Lovibond,
1995})), which consists of 42 items measuring symptoms of depression,
anxiety, and stress on a 4-point likert scale.

\paragraph{Brief Symptom Inventory.}\label{brief-symptom-inventory}

Overall psychopathological symptom burden was measured using the
\emph{Brief Symptom Inventory} (BSI; (\citeproc{ref-book}{Franke,
2002})), the short form of the Symptom Checklist-90-Revised (SCL-90-R;
Derogatis). The BSI consists of 53 items rated on a 5-point Likert scale
ranging from 0 (``not at all'') to 4 (``extremely''). Responses to 49
items are assigned to nine primary symptom dimensions, while four items
are evaluated separately. These symptom dimensions are summarized into
three global indices: the \emph{Global Severity Index} (GSI), reflecting
overall psychological distress; the \emph{Positive Symptom Distress
Index} (PSDI), indicating symptom intensity; and the \emph{Positive
Symptom Total} (PST), representing the number of reported symptoms.

\paragraph{Positive Mental Health
Scale.}\label{positive-mental-health-scale}

Positive mental health (PMH) was assessed with the nine-item \emph{PMH
scale} (\citeproc{ref-lukat2016}{Lukat et al., 2016}). Responses are
given on a 4-point Likert scale from 0 (disagree) to 3 (agree). Item
scores are summed to yield a total score ranging from 0 to 27, with
higher scores reflecting greater PMH. The scale has been validated as a
unidimensional measure with excellent internal consistency (Cronbach's α
= .93), good test--retest reliability (Pearson r = .74--.81), and
evidence of scalar invariance across samples and over time
(\citeproc{ref-lukat2016}{Lukat et al., 2016}). Furthermore, it shows
strong convergent and discriminant validity and is sensitive to
therapeutic change across diverse populations
(\citeproc{ref-lukat2016}{Lukat et al., 2016}).

\paragraph{Childhood Trauma
Questionnaire.}\label{childhood-trauma-questionnaire}

Early adverse experiences were assessed using the \emph{Childhood Trauma
Questionnaire} (CTQ) (\citeproc{ref-bernstein2003}{Bernstein et al.,
2003}), a widely used self-report instrument for the retrospective
assessment of childhood maltreatment. The CTQ measures five domains of
adverse experiences: emotional abuse, physical abuse, sexual abuse,
emotional neglect, and physical neglect. Items are rated on a 5-point
Likert scale ranging from 1 (``not at all'') to 5 (``very often''), with
higher scores indicating greater exposure to maltreatment. The German
version of the CTQ has demonstrated good psychometric properties,
including satisfactory reliability and validity in clinical samples
(\citeproc{ref-wingenfeld2010}{Wingenfeld et al., 2010}).

\paragraph{Clinical Global
Impression.}\label{clinical-global-impression}

Clinician-rated symptom severity and improvement were assessed using the
\emph{Clinical Global Impression} (CGI) scales.

\subparagraph{Severity Scale.}\label{severity-scale}

The \emph{CGI-Severity} scale evaluates the clinician's global
impression of the patient's current level of mental illness, based on
their total clinical experience with this population. The item asks:
\emph{``Considering your total clinical experience with this particular
population, how mentally ill is the patient at this time?''}

\subparagraph{Improvement Scale.}\label{improvement-scale}

Treatment-related change was assessed using the \emph{CGI-Improvement}
scale. Both patients and therapists rated overall improvement relative
to the beginning of therapy, regardless of whether the change was
attributed entirely to treatment. Patient and therapist versions differ
only in perspective but use equivalent response formats.

\subsubsection{Global Improvement}\label{global-improvement}

Global therapy outcome was assessed using a six-point global success
rating based on two items measuring perceived benefit and satisfaction
with therapy (\citeproc{ref-michalak2003}{Michalak et al., 2003}). These
items were completed by both patients and therapists and capture a
retrospective evaluation of treatment success. The items assess (1) the
extent to which expectations toward therapy have been fulfilled and (2)
the overall perceived benefit of therapy. Responses are given on a
6-point Likert scale ranging from 1 (``on the contrary / rather
harmful'') to 6 (``completely / very helpful'').

\subsubsection{Goal Attainment Scale}\label{goal-attainment-scale}

Individualized treatment outcomes were assessed using a goal attainment
measure inspired by the Goal Attainment Scaling approach
(\citeproc{ref-kiresuk1968}{Kiresuk \& Sherman, 1968}). At the beginning
of therapy, patients and therapists collaboratively define
individualized treatment goals. During the course and at the end of
therapy, patients and therapists retrospectively evaluated the extent to
which each of the predefined goals had been achieved.

Goal attainment was rated on a standardized six-point numerical scale
ranging from deterioration relative to the initial goal state (-1 =
moved away from the goal) to full goal attainment (4 = goal achieved).
The scale reflects patients' subjective assessment of goal progress,
with intermediate categories indicating partial progress toward the
respective goal.

For each patient, an overall goal attainment score was computed as the
mean rating across all individually defined goals, representing the
average subjective level of goal progress at the end of therapy.

\subsection{Measurement time points}\label{measurement-time-points}

\begin{longtable}[]{@{}
  >{\raggedright\arraybackslash}p{(\linewidth - 16\tabcolsep) * \real{0.1111}}
  >{\raggedright\arraybackslash}p{(\linewidth - 16\tabcolsep) * \real{0.1111}}
  >{\raggedright\arraybackslash}p{(\linewidth - 16\tabcolsep) * \real{0.1111}}
  >{\raggedright\arraybackslash}p{(\linewidth - 16\tabcolsep) * \real{0.1111}}
  >{\raggedright\arraybackslash}p{(\linewidth - 16\tabcolsep) * \real{0.1111}}
  >{\raggedright\arraybackslash}p{(\linewidth - 16\tabcolsep) * \real{0.1111}}
  >{\raggedright\arraybackslash}p{(\linewidth - 16\tabcolsep) * \real{0.1111}}
  >{\raggedright\arraybackslash}p{(\linewidth - 16\tabcolsep) * \real{0.1111}}
  >{\raggedright\arraybackslash}p{(\linewidth - 16\tabcolsep) * \real{0.1111}}@{}}
\toprule\noalign{}
\begin{minipage}[b]{\linewidth}\raggedright
Timepoint
\end{minipage} & \begin{minipage}[b]{\linewidth}\raggedright
DU-DI
\end{minipage} & \begin{minipage}[b]{\linewidth}\raggedright
DU-Prä
\end{minipage} & \begin{minipage}[b]{\linewidth}\raggedright
KZT1-DU4
\end{minipage} & \begin{minipage}[b]{\linewidth}\raggedright
KZT1-DUPost
\end{minipage} & \begin{minipage}[b]{\linewidth}\raggedright
KZT2-DUPost
\end{minipage} & \begin{minipage}[b]{\linewidth}\raggedright
LZT1-DUPost
\end{minipage} & \begin{minipage}[b]{\linewidth}\raggedright
LZT2-DUPost
\end{minipage} & \begin{minipage}[b]{\linewidth}\raggedright
Kat6
\end{minipage} \\
\midrule\noalign{}
\endhead
\bottomrule\noalign{}
\endlastfoot
\textbf{Explanation} & Pre-therapy, 4th contact & Pre-therapy, 6th
contact & 4th therapy session & 12th therapy session & 24th therapy
session & 45th therapy session & 60th therapy session & 6 month after
therapy \\
\textbf{Diagnosis} & X & & & & & & & \\
\textbf{Demographics} & & X & & & & & & \\
\textbf{BSI} & X & & X & X & X & X & X & X \\
\textbf{BDI-II} & & X & X & X & X & X & X & X \\
\textbf{DASS-42} & & X & X & X & X & X & X & X \\
\textbf{PMH} & & X & X & X & X & X & X & X \\
\textbf{CTQ} & & X & & & & & & \\
\textbf{CGI-S} & & X & & & & & & \\
\textbf{CGI-I} & & & & X & X & X & X & X \\
\textbf{Glob-Pt} & & & & X & X & X & X & X \\
\textbf{GAS} & & & & X & X & X & X & X \\
\end{longtable}

\subsection{Preprocessing}\label{preprocessing}

To streamline data collection, an automated transcription pipeline was
implemented: The handwritten text data is first recorded by trained
employees of the FBZ adult outpatient clinic using a mobile audio
recording device. Identifying features (e.g.~names, dates of birth,
location details) were replaced by placeholders during recording. The
transcription was carried out on local hardware using the open source
tool Whisper Large v2 (\url{https://github.com/openai/whisper}), a
state-of-the-art speech-to-text model
(\citeproc{ref-Radford2022}{Radford et al., 2022}). Each recording
begins with a structured introduction, including a patient
identification code, followed by responses to predefined questions. The
transcription pipeline automatically processes all audio recordings,
extracts the patient codes, and identifies responses to key questions.

As an additional data correction step, the exported transcription table
was screened for incomplete entries. Records with missing patient
identification codes or without any extracted text were automatically
flagged, exported for manual correction, and subsequently re-imported
and merged back into the original dataset. The corrected dataset was
then used for downstream analyses.

\subsection{Shared analytic framework}\label{shared-analytic-framework}

All substudies in this project are based on the same clinical cohort and
share a common analytic foundation. Specific operationalizations of text
inputs, outcomes, and analytic models differ across substudies and are
specified in the corresponding substudy sections below.

\textbf{Data scope and unit of analysis}. Across all substudies,
analyses focus on pre-therapy patient narratives collected during
intake. Language-based analyses are temporally ordered such that patient
narratives precede all clinical outcomes of interest. Outcomes may be
assessed cross-sectionally (at intake) or longitudinally (during or
after therapy), depending on the substudy. The unit of analysis is the
individual patient. Textual data collected during therapy are not used
as predictors in any analysis.

\textbf{Descriptive characterization of open-ended responses.} As a
shared descriptive foundation, responses to each open-ended prompt are
characterized with respect to engagement and heterogeneity. Descriptive
statistics include response length and an entropy-based lexical
diversity index (\citeproc{ref-shannon1948mathematical}{Shannon, 1948}),
computed across pooled responses per question. This characterization
provides a common empirical basis for interpreting prompt-specific
response patterns across substudies.

Because the open-ended questions were administered in two questionnaire
blocks, analyses are restricted to cases in which the respective
questionnaire was present. Within these blocks, item-level nonresponse
is summarized descriptively and interpreted as potentially informative
of selective responding.

\section{Preliminary results}\label{preliminary-results}

\subsection{Descriptive statistics}\label{descriptive-statistics}

\subsubsection{Demographics and context
factors}\label{demographics-and-context-factors}

\begin{longtable}[]{@{}
  >{\raggedright\arraybackslash}p{(\linewidth - 10\tabcolsep) * \real{0.2248}}
  >{\centering\arraybackslash}p{(\linewidth - 10\tabcolsep) * \real{0.0698}}
  >{\centering\arraybackslash}p{(\linewidth - 10\tabcolsep) * \real{0.0853}}
  >{\raggedright\arraybackslash}p{(\linewidth - 10\tabcolsep) * \real{0.4031}}
  >{\centering\arraybackslash}p{(\linewidth - 10\tabcolsep) * \real{0.1163}}
  >{\centering\arraybackslash}p{(\linewidth - 10\tabcolsep) * \real{0.1008}}@{}}
\caption{Table X. Descriptive statistics of patient demographics and
context factors.}\tabularnewline
\toprule\noalign{}
\begin{minipage}[b]{\linewidth}\raggedright
Variable
\end{minipage} & \begin{minipage}[b]{\linewidth}\centering
Valid n
\end{minipage} & \begin{minipage}[b]{\linewidth}\centering
Missing n
\end{minipage} & \begin{minipage}[b]{\linewidth}\raggedright
Level
\end{minipage} & \begin{minipage}[b]{\linewidth}\centering
Mean ± SD
\end{minipage} & \begin{minipage}[b]{\linewidth}\centering
n (\%)
\end{minipage} \\
\midrule\noalign{}
\endfirsthead
\toprule\noalign{}
\begin{minipage}[b]{\linewidth}\raggedright
Variable
\end{minipage} & \begin{minipage}[b]{\linewidth}\centering
Valid n
\end{minipage} & \begin{minipage}[b]{\linewidth}\centering
Missing n
\end{minipage} & \begin{minipage}[b]{\linewidth}\raggedright
Level
\end{minipage} & \begin{minipage}[b]{\linewidth}\centering
Mean ± SD
\end{minipage} & \begin{minipage}[b]{\linewidth}\centering
n (\%)
\end{minipage} \\
\midrule\noalign{}
\endhead
\bottomrule\noalign{}
\endlastfoot
\textbf{Age at therapy start} & 675 & 256 & & 41.48 ± 14.23 & \\
\textbf{Sex} & 675 & 256 & & & \\
& & & male & & 269 (39.9\%) \\
& & & female & & 406 (60.1\%) \\
\textbf{In relationship} & 381 & 550 & & & \\
& & & yes & & 220 (57.7\%) \\
& & & no & & 161 (42.3\%) \\
\textbf{Marital status} & 381 & 550 & & & \\
& & & single & & 193 (50.7\%) \\
& & & married & & 107 (28.1\%) \\
& & & divorced & & 52 (13.6\%) \\
& & & seperated & & 16 (4.2\%) \\
& & & widowed & & 5 (1.3\%) \\
& & & other & & 8 (2.1\%) \\
\textbf{General education} & 381 & 550 & & & \\
& & & other & & 5 (1.3\%) \\
& & & student & & 4 (1.0\%) \\
& & & no school-leaving certificate & & 9 (2.4\%) \\
& & & lower secondary school certificate & & 66 (17.3\%) \\
& & & intermediate secondary school certificate & & 114 (29.9\%) \\
& & & higher education entrance qualification & & 183 (48.0\%) \\
\textbf{Vocational qualification} & 381 & 550 & & & \\
& & & Currently in vocational training or studying & & 40 (10.5\%) \\
& & & No vocational qualification & & 45 (11.8\%) \\
& & & Apprenticeship / vocational training & & 217 (57.0\%) \\
& & & University or university of applied sciences degree & & 54
(14.2\%) \\
& & & Other & & 25 (6.6\%) \\
\textbf{Work ability status} & 381 & 550 & & & \\
& & & Other & & 30 (7.9\%) \\
& & & Able to work & & 208 (54.6\%) \\
& & & Unable to work (on sick leave) & & 118 (31.0\%) \\
& & & Disability pension & & 15 (3.9\%) \\
& & & Old-age pension & & 10 (2.6\%) \\
\textbf{Previous psychotherapy} & 582 & 349 & & & \\
& & & no prior treatment & & 179 (30.8\%) \\
& & & outpatient psychotherapy & & 100 (17.2\%) \\
& & & inpatient psychotherapy & & 154 (26.5\%) \\
& & & both & & 134 (23.0\%) \\
& & & exact specification not available & & 15 (2.6\%) \\
\textbf{CGI severity} & 582 & 349 & & & \\
& & & Not assessable & & 4 (0.7\%) \\
& & & Normal, not at all ill & & 2 (0.3\%) \\
& & & Borderline mentally ill & & 9 (1.5\%) \\
& & & Mildly ill & & 28 (4.8\%) \\
& & & Moderately ill & & 162 (27.8\%) \\
& & & Markedly ill & & 287 (49.3\%) \\
& & & Severely ill & & 87 (14.9\%) \\
& & & Among the most extremely ill patients & & 3 (0.5\%) \\
\end{longtable}

\begin{longtable}[]{@{}
  >{\raggedright\arraybackslash}p{(\linewidth - 10\tabcolsep) * \real{0.3165}}
  >{\centering\arraybackslash}p{(\linewidth - 10\tabcolsep) * \real{0.1139}}
  >{\centering\arraybackslash}p{(\linewidth - 10\tabcolsep) * \real{0.1392}}
  >{\raggedright\arraybackslash}p{(\linewidth - 10\tabcolsep) * \real{0.0886}}
  >{\centering\arraybackslash}p{(\linewidth - 10\tabcolsep) * \real{0.1772}}
  >{\centering\arraybackslash}p{(\linewidth - 10\tabcolsep) * \real{0.1646}}@{}}
\caption{Table X. Descriptive statistics of therapist demographics and
context factors.}\tabularnewline
\toprule\noalign{}
\begin{minipage}[b]{\linewidth}\raggedright
Variable
\end{minipage} & \begin{minipage}[b]{\linewidth}\centering
Valid n
\end{minipage} & \begin{minipage}[b]{\linewidth}\centering
Missing n
\end{minipage} & \begin{minipage}[b]{\linewidth}\raggedright
Level
\end{minipage} & \begin{minipage}[b]{\linewidth}\centering
Mean ± SD
\end{minipage} & \begin{minipage}[b]{\linewidth}\centering
n (\%)
\end{minipage} \\
\midrule\noalign{}
\endfirsthead
\toprule\noalign{}
\begin{minipage}[b]{\linewidth}\raggedright
Variable
\end{minipage} & \begin{minipage}[b]{\linewidth}\centering
Valid n
\end{minipage} & \begin{minipage}[b]{\linewidth}\centering
Missing n
\end{minipage} & \begin{minipage}[b]{\linewidth}\raggedright
Level
\end{minipage} & \begin{minipage}[b]{\linewidth}\centering
Mean ± SD
\end{minipage} & \begin{minipage}[b]{\linewidth}\centering
n (\%)
\end{minipage} \\
\midrule\noalign{}
\endhead
\bottomrule\noalign{}
\endlastfoot
\textbf{Age at therapy start} & 675 & 256 & & 28.35 ± 4.12 & \\
\textbf{Sex} & 675 & 256 & & & \\
& & & male & & 101 (15.0\%) \\
& & & female & & 574 (85.0\%) \\
\end{longtable}

\subsubsection{Textual data}\label{textual-data}

Descriptive analyses revealed substantial heterogeneity across questions
in response rates, length, and lexical diversity. Holistic questions
such as problem description (q10) and therapy goals (q12) showed low
missingness, longer median response lengths, and high lexical diversity,
indicating that patients readily produced extended and heterogeneous
narratives when asked to reflect broadly on their difficulties or
desired changes. In contrast, prompts targeting causal connections (q4)
or constraints (q11) frequently elicited short or missing responses and
exhibited lower diversity, consistent with more constrained or
confirmatory response formats. Notably, lexical diversity varied
independently of response length, suggesting that some prompts elicited
shared narrative scripts despite moderate verbosity (e.g., q7 Expected
Improvements).

\begin{longtable}[]{@{}
  >{\raggedright\arraybackslash}p{(\linewidth - 14\tabcolsep) * \real{0.1947}}
  >{\centering\arraybackslash}p{(\linewidth - 14\tabcolsep) * \real{0.0796}}
  >{\centering\arraybackslash}p{(\linewidth - 14\tabcolsep) * \real{0.1150}}
  >{\centering\arraybackslash}p{(\linewidth - 14\tabcolsep) * \real{0.1416}}
  >{\centering\arraybackslash}p{(\linewidth - 14\tabcolsep) * \real{0.1239}}
  >{\centering\arraybackslash}p{(\linewidth - 14\tabcolsep) * \real{0.0973}}
  >{\centering\arraybackslash}p{(\linewidth - 14\tabcolsep) * \real{0.0973}}
  >{\centering\arraybackslash}p{(\linewidth - 14\tabcolsep) * \real{0.1504}}@{}}
\caption{Table X. Descriptive statistics of open-ended responses by
question.}\tabularnewline
\toprule\noalign{}
\begin{minipage}[b]{\linewidth}\raggedright
Question
\end{minipage} & \begin{minipage}[b]{\linewidth}\centering
n\_total
\end{minipage} & \begin{minipage}[b]{\linewidth}\centering
missing\_pct
\end{minipage} & \begin{minipage}[b]{\linewidth}\centering
very\_short\_pct
\end{minipage} & \begin{minipage}[b]{\linewidth}\centering
median\_words
\end{minipage} & \begin{minipage}[b]{\linewidth}\centering
p10\_words
\end{minipage} & \begin{minipage}[b]{\linewidth}\centering
p90\_words
\end{minipage} & \begin{minipage}[b]{\linewidth}\centering
diversity\_index
\end{minipage} \\
\midrule\noalign{}
\endfirsthead
\toprule\noalign{}
\begin{minipage}[b]{\linewidth}\raggedright
Question
\end{minipage} & \begin{minipage}[b]{\linewidth}\centering
n\_total
\end{minipage} & \begin{minipage}[b]{\linewidth}\centering
missing\_pct
\end{minipage} & \begin{minipage}[b]{\linewidth}\centering
very\_short\_pct
\end{minipage} & \begin{minipage}[b]{\linewidth}\centering
median\_words
\end{minipage} & \begin{minipage}[b]{\linewidth}\centering
p10\_words
\end{minipage} & \begin{minipage}[b]{\linewidth}\centering
p90\_words
\end{minipage} & \begin{minipage}[b]{\linewidth}\centering
diversity\_index
\end{minipage} \\
\midrule\noalign{}
\endhead
\bottomrule\noalign{}
\endlastfoot
Problem development & 931 & 47.046187 & 2.577873 & 18 & 6 & 50.0 &
731.0489 \\
Extra stressors & 931 & 49.838883 & 9.129968 & 10 & 2 & 36.0 &
636.8977 \\
Pre-onset changes & 931 & 43.179377 & 15.037594 & 9 & 1 & 37.0 &
555.4025 \\
Event connection & 931 & 47.046187 & 23.630505 & 5 & 1 & 28.8 &
477.7177 \\
Physical symptoms & 931 & 46.401719 & 16.433942 & 6 & 2 & 19.0 &
500.2739 \\
Problem causes & 931 & 49.087003 & 12.137487 & 8 & 2 & 30.0 &
587.9515 \\
Expected improvements & 931 & 44.683136 & 5.477981 & 11 & 4 & 31.0 &
377.6039 \\
Environment response & 931 & 43.286788 & 12.137487 & 10 & 2 & 32.3 &
363.8619 \\
No change required & 931 & 71.321160 & 7.626208 & 6 & 2 & 19.0 &
275.3418 \\
Problem description & 931 & 11.170784 & 4.511278 & 22 & 6 & 50.4 &
687.7989 \\
Impacted life areas & 931 & 12.996778 & 31.256713 & 6 & 2 & 31.1 &
438.1475 \\
Therapy goals & 931 & 9.881847 & 5.477981 & 15 & 5 & 36.0 & 435.7618 \\
\end{longtable}

\subsubsection{Diagnosis}\label{diagnosis-1}

\begin{longtable}[]{@{}
  >{\raggedright\arraybackslash}p{(\linewidth - 4\tabcolsep) * \real{0.5800}}
  >{\centering\arraybackslash}p{(\linewidth - 4\tabcolsep) * \real{0.1900}}
  >{\centering\arraybackslash}p{(\linewidth - 4\tabcolsep) * \real{0.2300}}@{}}
\caption{DSM-5 diagnoses in the sample (absolute
frequencies)}\tabularnewline
\toprule\noalign{}
\begin{minipage}[b]{\linewidth}\raggedright
DSM-5 diagnosis
\end{minipage} & \begin{minipage}[b]{\linewidth}\centering
All diagnoses (n)
\end{minipage} & \begin{minipage}[b]{\linewidth}\centering
Primary diagnoses (n)
\end{minipage} \\
\midrule\noalign{}
\endfirsthead
\toprule\noalign{}
\begin{minipage}[b]{\linewidth}\raggedright
DSM-5 diagnosis
\end{minipage} & \begin{minipage}[b]{\linewidth}\centering
All diagnoses (n)
\end{minipage} & \begin{minipage}[b]{\linewidth}\centering
Primary diagnoses (n)
\end{minipage} \\
\midrule\noalign{}
\endhead
\bottomrule\noalign{}
\endlastfoot
Major Depression (all subtypes) & 331 & 251 \\
Persistent Depressive Disorder (Dysthymia) & 124 & 89 \\
Social Anxiety Disorder (Social Phobia) & 103 & 56 \\
Posttraumatic Stress Disorder & 65 & 33 \\
Agoraphobia & 59 & 34 \\
Panic Disorder & 54 & 28 \\
Adjustment Disorder with Depressed Mood & 40 & 40 \\
Generalized Anxiety Disorder & 36 & 22 \\
Borderline Personality Disorder & 31 & 20 \\
Somatic Symptom Disorder & 28 & 20 \\
Obsessive-Compulsive Disorder & 25 & 16 \\
Adjustment Disorder with Mixed Anxiety and Depressed Mood & 16 & 16 \\
Specific Phobia, Situational Type & 15 & 6 \\
Binge-Eating Disorder & 14 & 3 \\
Insomnia Disorder & 14 & 3 \\
Separation Anxiety Disorder & 13 & 3 \\
Adjustment Disorder, Unspecified & 11 & 9 \\
Illness Anxiety Disorder & 11 & 4 \\
Alcohol Use Disorder, Moderate & 10 & 3 \\
Specific Phobia, Natural Environment Type & 10 & 1 \\
Other diagnoses (\textless10) & 166 & 75 \\
\end{longtable}

Diagnoses were collapsed for display. Only diagnoses with ≥10
occurrences are shown individually; remaining diagnoses are summarized
as ``Other diagnoses (\textless10)''.

\subsubsection{psychometric measures}\label{psychometric-measures-1}

Due to substantial dropout and differences in treatment duration,
descriptive statistics across therapy phases are based on changing
subsamples. In particular, patients continuing into long-term therapy
represent a more severe and chronic subgroup. To avoid conflating
symptom severity with symptom change, two complementary descriptive
tables are reported.

Table 1 presents descriptive statistics for symptom measures at each
assessment timepoint. Sample sizes decrease substantially across therapy
phases, reflecting treatment completion and dropout. Consequently, means
at later timepoints do not represent longitudinal change within
individuals and should not be interpreted as symptom worsening.

\begin{longtable*}{lcccccccc}
\toprule
Scale & DI & DU-Prä & KZT1-DU4 & KZT1-DUPost & KZT2-DUPost & LZT1-DUPost & LZT2-DUPost & DU-Kat6\\
\midrule
\textbf{BDI} &  &  &  &  &  &  &  & \\
\hspace{1em}Sum &  & 23.45 ± 12.52\\ n = 595 & 19.12 ± 11.23\\ n = 272 & 16.42 ± 12.16\\ n = 390 & 14.98 ± 12.40\\ n = 299 & 14.53 ± 12.35\\ n = 118 & 15.34 ± 12.42\\ n = 29 & 12.96 ± 11.53\\ n = 211\\
\textbf{BSI} &  &  &  &  &  &  &  & \\
\hspace{1em}GSI & 1.24 ± 0.71\\ n = 495 &  & 1.08 ± 0.67\\ n = 275 & 1.01 ± 0.68\\ n = 403 & 0.84 ± 0.67\\ n = 314 & 1.03 ± 0.73\\ n = 124 & 1.07 ± 0.62\\ n = 30 & 0.84 ± 0.68\\ n = 231\\
\textbf{CGI} &  &  &  &  &  &  &  & \\
\addlinespace
\hspace{1em}improvement patient &  &  &  & 2.36 ± 1.13\\ n = 358 & 2.17 ± 1.08\\ n = 272 & 1.90 ± 0.93\\ n = 117 & 1.77 ± 0.57\\ n = 30 & 2.26 ± 1.41\\ n = 220\\
\hspace{1em}improvement therapist &  &  &  & 2.73 ± 0.83\\ n = 269 & 2.33 ± 0.93\\ n = 182 & 2.03 ± 0.65\\ n = 68 & 1.89 ± 0.58\\ n = 18 & 1.74 ± 1.29\\ n = 53\\
\textbf{DASS} &  &  &  &  &  &  &  & \\
\hspace{1em}Anxiety &  & 11.62 ± 8.27\\ n = 425 & 9.16 ± 7.30\\ n = 146 & 8.07 ± 7.78\\ n = 283 & 7.22 ± 7.19\\ n = 236 & 8.26 ± 7.93\\ n = 85 & 6.81 ± 8.11\\ n = 21 & 6.66 ± 7.17\\ n = 223\\
\hspace{1em}Depression &  & 19.37 ± 11.10\\ n = 425 & 16.38 ± 10.18\\ n = 146 & 13.03 ± 10.60\\ n = 283 & 12.60 ± 11.09\\ n = 236 & 11.13 ± 9.92\\ n = 85 & 11.95 ± 9.54\\ n = 21 & 11.39 ± 10.70\\ n = 223\\
\addlinespace
\hspace{1em}Stress &  & 18.91 ± 9.29\\ n = 425 & 16.23 ± 8.74\\ n = 146 & 14.56 ± 9.31\\ n = 283 & 12.56 ± 9.33\\ n = 236 & 13.73 ± 9.83\\ n = 85 & 12.90 ± 9.04\\ n = 21 & 12.66 ± 9.97\\ n = 223\\
\hspace{1em}Total &  & 49.90 ± 24.23\\ n = 425 & 41.77 ± 22.44\\ n = 146 & 35.65 ± 24.47\\ n = 283 & 32.38 ± 24.56\\ n = 236 & 33.12 ± 25.07\\ n = 85 & 31.67 ± 23.42\\ n = 21 & 30.72 ± 25.00\\ n = 223\\
\textbf{GAS} &  &  &  &  &  &  &  & \\
\hspace{1em}pt mean &  &  &  & 1.48 ± 1.03\\ n = 282 & 2.03 ± 1.10\\ n = 235 & 2.12 ± 1.08\\ n = 101 & 2.48 ± 0.84\\ n = 24 & 2.25 ± 1.15\\ n = 102\\
\hspace{1em}th mean &  &  &  & 1.49 ± 1.02\\ n = 217 & 2.09 ± 1.08\\ n = 160 & 2.53 ± 0.94\\ n = 57 & 2.98 ± 0.62\\ n = 14 & \\
\addlinespace
\textbf{Global ratings} &  &  &  &  &  &  &  & \\
\hspace{1em}pt benefit &  &  &  & 4.59 ± 1.03\\ n = 404 & 4.95 ± 0.96\\ n = 311 & 5.07 ± 0.90\\ n = 123 & 5.40 ± 0.62\\ n = 30 & 4.72 ± 1.20\\ n = 230\\
\hspace{1em}pt satisfaction &  &  &  & 4.15 ± 1.06\\ n = 404 & 4.48 ± 1.08\\ n = 311 & 4.50 ± 0.98\\ n = 123 & 4.70 ± 0.75\\ n = 30 & 4.33 ± 1.19\\ n = 230\\
\hspace{1em}th benefit &  &  &  & 3.89 ± 0.94\\ n = 311 & 4.36 ± 1.00\\ n = 210 & 4.76 ± 0.80\\ n = 72 & 5.17 ± 0.51\\ n = 18 & \\
\hspace{1em}th satisfaction &  &  &  & 3.67 ± 0.96\\ n = 311 & 4.09 ± 1.09\\ n = 210 & 4.44 ± 0.99\\ n = 72 & 4.56 ± 0.78\\ n = 18 & \\
\addlinespace
\textbf{PMH} &  &  &  &  &  &  &  & \\
\hspace{1em}Sum &  & 10.64 ± 5.79\\ n = 588 & 11.50 ± 6.00\\ n = 274 & 13.03 ± 6.43\\ n = 401 & 14.05 ± 6.53\\ n = 316 & 13.44 ± 6.24\\ n = 121 & 12.70 ± 5.83\\ n = 30 & 15.24 ± 6.79\\ n = 230\\
\bottomrule
\end{longtable*}

Table 2 reports descriptive statistics restricted to patients with
complete data at both pre-treatment and post-treatment assessments. For
complete-case descriptive analyses, the post-treatment endpoint was
defined as the assessment at the end of the second short-term therapy
phase (KZT2-DUPost), as this represents the last common assessment
across the majority of patients. This table provides a descriptive
approximation of within-person change, independent of selective dropout.

\begin{longtable*}{lcccccc}
\toprule
Scale & Baseline & Endpoint & Baseline M ± SD & Endpoint M ± SD & ΔM (End − Base) & n (complete)\\
\midrule
\textbf{BDI} &  &  &  &  &  & \\
\hspace{1em}Sum & DU-Prä & KZT2-DUPost & 22.63 ± 11.98 & 14.98 ± 12.40 & -7.65 & 299\\
\textbf{BSI} &  &  &  &  &  & \\
\hspace{1em}GSI & DI & KZT2-DUPost & 1.19 ± 0.67 & 0.83 ± 0.66 & -0.36 & 253\\
\textbf{DASS} &  &  &  &  &  & \\
\addlinespace
\hspace{1em}Anxiety & DU-Prä & KZT2-DUPost & 11.38 ± 7.99 & 6.94 ± 6.95 & -4.43 & 192\\
\hspace{1em}Depression & DU-Prä & KZT2-DUPost & 19.77 ± 11.41 & 12.82 ± 11.47 & -6.95 & 192\\
\hspace{1em}Stress & DU-Prä & KZT2-DUPost & 19.02 ± 9.69 & 12.28 ± 9.22 & -6.74 & 192\\
\hspace{1em}Total & DU-Prä & KZT2-DUPost & 50.16 ± 24.37 & 32.04 ± 24.61 & -18.12 & 192\\
\textbf{PMH} &  &  &  &  &  & \\
\addlinespace
\hspace{1em}Sum & DU-Prä & KZT2-DUPost & 10.60 ± 5.70 & 13.94 ± 6.56 & 3.34 & 297\\
\bottomrule
\end{longtable*}

\section{Research areas}\label{research-areas}

This project comprises multiple complementary research areas that share
a common dataset and preprocessing pipeline but address distinct
research questions. To avoid analytic flexibility and to maximize
interpretability, each substudy pre-specifies its text inputs, outcomes,
and evaluation strategy.

\subsection{Research area 1: Language as clinical assessment
(cross-sectional)}\label{research-area-1-language-as-clinical-assessment-cross-sectional}

\textbf{Core research question.} Does patients' pre-therapy language
provide clinically meaningful information about symptom severity and
functioning, and does it add predictive value beyond established
self-report questionnaires?

\textbf{Text inputs.} Q1 Problem developtment, Q10 \emph{Problem
description,} Q12 \emph{Therapy goals}

\textbf{Outcomes.}

\begin{itemize}
\item
  Convergent validity outcomes: Cross-sectional self-report symptom
  (BDI-II, BSI/GSI, DASS-42) and well-being (PMH) measures
\item
  Incremental validity outcomes: Clinician-rated symptom severity
  (CGI-Severity).
\end{itemize}

\textbf{Models \& evaluation.} Convergent validity with self-report
symptom measures (BDI-II, BSI/GSI, DASS-42) will be evaluated by
estimating the proportion of symptom variance explained by
language-based representations. Incremental validity will be tested by
assessing whether language explains additional variance in
clinician-rated symptom severity (CGI-Severity) beyond self-report
questionnaires.

The Sequential Evaluation with Model Pre-registration
(\citeproc{ref-kjell_ganesan_boyd_oltmanns_rivero_feltman_carr_luft_kotov_schwartz_2024}{Kjell,
Ganesan, et al., 2024}) framework will be implemented to ensure robust
model development, mitigating overfitting and enabling unbiased
performance evaluation. Additionally, evaluating models on prospective
data will simulate real-world clinical deployment by assessing
performance on new, unseen patient data.

Contextual embeddings derived from pretrained LLMs will be linked to
clinical outcomes using state-of-the-art prediction models, including
ridge regression (\citeproc{ref-Hoerl1970}{Hoerl \& Kennard, 1970}),
lasso regression
(\citeproc{ref-10.1111ux2fj.2517-6161.1996.tb02080.x}{Tibshirani,
1996}), and random forests (\citeproc{ref-598994}{Ho, 1995}).

\subsection{Research area 2: Patient narratives and
themes}\label{research-area-2-patient-narratives-and-themes}

\textbf{Core research question.} How do patients conceptualize the their
mental health problems, as well as anticipated improvements, and how do
these narrative patterns vary across diagnostic categories, baseline
severity or sociodemographic (e.g.~age, gender) and anamnestic
(e.g.~childhood trauma) characteristics.

\textbf{Text inputs.} Q1 \emph{Problem development}, Q2 \emph{Extra
stressors}, Q3 \emph{Pre-onset changes}, Q6 \emph{Problem causes}, Q7
\emph{Expected improvements}, Q8 \emph{Environment response}.
Question-wise modeling.

\textbf{Outcomes.} No single predictive outcome is specified, as the
primary aim is interpretive and theory-generating rather than
predictive.

\textbf{Models \& evaluation.} Topic modeling; interpretive labeling;
comparison of topic prevalence across patient groups.

\subsection{Research area 3: Predicting treatment response and goal
attainment (longitudinal
prediction)}\label{research-area-3-predicting-treatment-response-and-goal-attainment-longitudinal-prediction}

\textbf{Core research question.} Can patients' pre-therapy language
predict clinically meaningful treatment outcomes beyond baseline symptom
severity and demographic or anamnestic characteristics?

\textbf{Conceptual contribution.} This substudy extends language-based
mental health assessment from cross-sectional validity to prognostic
utility. While prior work has primarily examined whether language
reflects current symptom severity, less is known about whether
pre-therapy narratives contain information relevant for predicting
future treatment response.

\textbf{Text inputs.} Primary: Q10, Q12, Q7, Q1; Secondary:
concatenation models.

\textbf{Outcomes.}

\begin{itemize}
\item
  Clinician-rated improvement (CGI-Improvement) at 6 months-follow-up
  assessment.
\item
  Global therapy outcome ratings (patient- and therapist-reported) at 6
  months-follow-up assessment.
\item
  Goal Attainment Scale at 6 months-follow-up assessment.
\item
  Pre-Post differences in symptom (BDI-II, BSI/GSI, DASS-42) and
  well-being (PMH) scales .
\end{itemize}

\textbf{Models \& evaluation.}

\begin{itemize}
\item
  Baseline prognostic models: Demographic variables and baseline symptom
  severity and well-being (e.g., BDI-II, BSI/GSI, PMH, CGI-Severity).
\item
  Language-augmented models: Baseline predictors plus pre-therapy
  language representations.
\end{itemize}

\subsection{Research area 4: Prompt-based LLMs for interpretable
language
assessment}\label{research-area-4-prompt-based-llms-for-interpretable-language-assessment}

\textbf{Core research question.} Can prompt-based LLMs reliably extract
interpretable, clinically meaningful scores from pre-therapy narratives,
and do these language measures add incremental value beyond (a)
standardized questionnaires and (b) embedding-based language
representations?

\section{References}

\phantomsection\label{refs}
\begin{CSLReferences}{1}{0}
\bibitem[\citeproctext]{ref-Beck1996}
Beck, A. T., Steer, R. A., \& Brown, G. K. (1996). \emph{{BDI-II}: {Beck
Depression Inventory} manual} (2nd ed.). Psychological Corporation.

\bibitem[\citeproctext]{ref-bernstein2003}
Bernstein, D. P., Stein, J. A., Newcomb, M. D., Walker, E., Pogge, D.,
Ahluvalia, T., Stokes, J., Handelsman, L., Medrano, M., Desmond, D., \&
Zule, W. (2003). Development and validation of a brief screening version
of the Childhood Trauma Questionnaire. \emph{Child Abuse \& Neglect},
\emph{27}(2), 169--190.
\url{https://doi.org/10.1016/s0145-2134(02)00541-0}

\bibitem[\citeproctext]{ref-chorpita2011evidence}
Chorpita, B. F., Daleiden, E. L., Ebesutani, C., Young, J., Becker, K.
D., Nakamura, B. J., Phillips, L., Ward, A., Lynch, R., Trent, L., et
al. (2011). Evidence-based treatments for children and adolescents: {An}
updated review of indicators of efficacy and effectiveness.
\emph{Clinical Psychology: Science and Practice}, \emph{18}(2),
154--172.

\bibitem[\citeproctext]{ref-devlin-etal-2019-bert}
Devlin, J., Chang, M.-W., Lee, K., \& Toutanova, K. (2019). {BERT}:
{Pre-training} of deep bidirectional transformers for language
understanding. In J. Burstein, C. Doran, \& T. Solorio (Eds.),
\emph{Proceedings of the 2019 conference of the north {American} chapter
of the association for computational linguistics: {Human} language
technologies, volume 1 (long and short papers)} (pp. 4171--4186).
Association for Computational Linguistics.
\url{https://doi.org/10.18653/v1/N19-1423}

\bibitem[\citeproctext]{ref-eberhardt2024}
Eberhardt, S. T., Schaffrath, J., Moggia, D., Schwartz, B., Jaehde, M.,
Rubel, J. A., Baur, T., André, E., \& Lutz, W. (2024). Decoding
emotions: Exploring the validity of sentiment analysis in psychotherapy.
\emph{Psychotherapy Research}, \emph{35}(2), 174--189.
\url{https://doi.org/10.1080/10503307.2024.2322522}

\bibitem[\citeproctext]{ref-eichstaedt2018}
Eichstaedt, J. C., Smith, R. J., Merchant, R. M., Ungar, L. H.,
Crutchley, P., Preoţiuc-Pietro, D., Asch, D. A., \& Schwartz, H. A.
(2018). Facebook language predicts depression in medical records.
\emph{Proceedings of the National Academy of Sciences}, \emph{115}(44),
11203--11208. \url{https://doi.org/10.1073/pnas.1802331115}

\bibitem[\citeproctext]{ref-book}
Franke, G. (2002). \emph{Franke, g.h. (2000). BSI. Brief symptom
inventory - deutsche version. Manual. Göttingen: beltz.}

\bibitem[\citeproctext]{ref-598994}
Ho, T. K. (1995). Random decision forests. \emph{Proceedings of 3rd
International Conference on Document Analysis and Recognition},
\emph{1}, 278--282 vol.1.
\url{https://doi.org/10.1109/ICDAR.1995.598994}

\bibitem[\citeproctext]{ref-Hoerl1970}
Hoerl, A. E., \& Kennard, R. W. (1970). Ridge {Regression}: {Biased
Estimation} for {Nonorthogonal Problems}. \emph{Technometrics},
\emph{12}(1), 55--67.
\url{https://doi.org/10.1080/00401706.1970.10488634}

\bibitem[\citeproctext]{ref-Jensen-Doss2008}
Jensen-Doss, A., \& Weisz, J. R. (2008). Diagnostic agreement predicts
treatment process and outcomes in youth mental health clinics.
\emph{Journal of Consulting and Clinical Psychology}, \emph{76}(5),
711--722. \url{https://doi.org/10.1037/0022-006X.76.5.711}

\bibitem[\citeproctext]{ref-kiresuk1968}
Kiresuk, T. J., \& Sherman, R. E. (1968). Goal attainment scaling: A
general method for evaluating comprehensive community mental health
programs. \emph{Community Mental Health Journal}, \emph{4}(6), 443--453.
\url{https://doi.org/10.1007/BF01530764}

\bibitem[\citeproctext]{ref-kjell_ganesan_boyd_oltmanns_rivero_feltman_carr_luft_kotov_schwartz_2024}
Kjell, O. N. E., Ganesan, A. V., Boyd, R., Oltmanns, J. R., Rivero, A.,
Feltman, S., Carr, M. A., Luft, B. J., Kotov, R., \& Schwartz, H. A.
(2024). \emph{Demonstrating high validity of a new {AI-language}
assessment of {PTSD}: A sequential evaluation with model
pre-registration}. PsyArXiv. \url{https://doi.org/10.31234/osf.io/xw24e}

\bibitem[\citeproctext]{ref-kjell2019semantic}
Kjell, O. N. E., Kjell, K., Garcia, D., \& Sikström, S. (2019). Semantic
measures: {Using} natural language processing to measure, differentiate,
and describe psychological constructs. \emph{Psychological Methods},
\emph{24}(1), 92--115. \url{https://doi.org/10.1037/met0000191}

\bibitem[\citeproctext]{ref-Kjell2024}
Kjell, O. N. E., Kjell, K., \& Schwartz, H. A. (2024). Beyond rating
scales: {With} targeted evaluation, large language models are poised for
psychological assessment. \emph{Psychiatry Research}, \emph{333},
115667. \url{https://doi.org/10.1016/j.psychres.2023.115667}

\bibitem[\citeproctext]{ref-Kjell2022}
Kjell, O. N. E., Sikström, S., Kjell, K., \& Schwartz, H. A. (2022).
Natural language analyzed with {AI-based} transformers predict
traditional subjective well-being measures approaching the theoretical
upper limits in accuracy. \emph{Scientific Reports}, \emph{12}(1), 3918.
\url{https://doi.org/10.1038/s41598-022-07520-w}

\bibitem[\citeproctext]{ref-Likert1932}
Likert, R. (1932). A technique for the measurement of attitudes.
\emph{Arch. Psychol.}, \emph{140}(55).

\bibitem[\citeproctext]{ref-Lovibond1995}
Lovibond, P. F., \& Lovibond, S. H. (1995). The structure of negative
emotional states: Comparison of the Depression Anxiety Stress Scales
(DASS) with the Beck Depression and Anxiety Inventories. \emph{Behaviour
Research and Therapy}, \emph{33}(3), 335--343.
\url{https://doi.org/10.1016/0005-7967(94)00075-u}

\bibitem[\citeproctext]{ref-lukat2016}
Lukat, J., Margraf, J., Lutz, R., Veld, W. M. van der, \& Becker, E. S.
(2016). Psychometric properties of the positive mental health scale
(PMH-scale). \emph{BMC Psychology}, \emph{4}(1), 8.
\url{https://doi.org/10.1186/s40359-016-0111-x}

\bibitem[\citeproctext]{ref-Lutz2022}
Lutz, W., Schwartz, B., \& Delgadillo, J. (2022). Measurement-{Based}
and {Data-Informed Psychological Therapy}. \emph{Annual Review of
Clinical Psychology}, \emph{18}(1), 71--98.
\url{https://doi.org/10.1146/annurev-clinpsy-071720-014821}

\bibitem[\citeproctext]{ref-margraf2021}
Margraf, J., Cwik, J. C., Brachel, R. von, Suppiger, A., \& Schneider,
S. (2021). \emph{DIPS open access 1.2: Diagnostisches interview bei
psychischen störungen}. \url{https://doi.org/10.46586/rub.172.149}

\bibitem[\citeproctext]{ref-michalak2003}
Michalak, J., Kosfelder, J., Meyer, F., \& Schulte, D. (2003). Messung
des Therapieerfolgs. \emph{Zeitschrift für Klinische Psychologie und
Psychotherapie}, \emph{32}(2), 94--103.
\url{https://doi.org/10.1026/0084-5345.32.2.94}

\bibitem[\citeproctext]{ref-Radford2022}
Radford, A., Kim, J. W., Xu, T., Brockman, G., McLeavey, C., \&
Sutskever, I. (2022). Robust {Speech Recognition} via {Large-Scale Weak
Supervision}. \emph{arXiv Preprint arXiv: 2212.04356}.
\url{https://arxiv.org/abs/2212.04356}

\bibitem[\citeproctext]{ref-shannon1948mathematical}
Shannon, C. E. (1948). A mathematical theory of communication. \emph{The
Bell System Technical Journal}, \emph{27}(3), 379423.

\bibitem[\citeproctext]{ref-tanana2021}
Tanana, M. J., Soma, C. S., Kuo, P. B., Bertagnolli, N. M., Dembe, A.,
Pace, B. T., Srikumar, V., Atkins, D. C., \& Imel, Z. E. (2021). How do
you feel? Using natural language processing to automatically rate
emotion in psychotherapy. \emph{Behavior Research Methods},
\emph{53}(5), 2069--2082.
\url{https://doi.org/10.3758/s13428-020-01531-z}

\bibitem[\citeproctext]{ref-10.1111ux2fj.2517-6161.1996.tb02080.x}
Tibshirani, R. (1996). Regression shrinkage and selection via the lasso.
\emph{Journal of the Royal Statistical Society: Series B
(Methodological)}, \emph{58}(1), 267--288.
\url{https://doi.org/10.1111/j.2517-6161.1996.tb02080.x}

\bibitem[\citeproctext]{ref-Vaswani2017}
Vaswani, A., Shazeer, N., Parmar, N., Uszkoreit, J., Jones, L., Gomez,
A. N., Kaiser, L., \& Polosukhin, I. (2017). Attention is {All} you
{Need}. \emph{Advances in {Neural Information Processing Systems}},
\emph{30}.

\bibitem[\citeproctext]{ref-wampold2015great}
Wampold, B. E., \& Imel, Z. E. (2015). \emph{The great psychotherapy
debate: {The} evidence for what makes psychotherapy work}. Routledge.

\bibitem[\citeproctext]{ref-wingenfeld2010}
Wingenfeld, K., Spitzer, C., Mensebach, C., Grabe, H., Hill, A., Gast,
U., Schlosser, N., Höpp, H., Beblo, T., \& Driessen, M. (2010). Die
deutsche Version des Childhood Trauma Questionnaire (CTQ): Erste Befunde
zu den psychometrischen Kennwerten. \emph{PPmP - Psychotherapie ·
Psychosomatik · Medizinische Psychologie}, \emph{60}(08), e13--e13.
\url{https://doi.org/10.1055/s-0030-1253494}

\bibitem[\citeproctext]{ref-WorldHealthOrganization2017}
World Health Organization. (2017). \emph{Depression and other common
mental disorders: {Global} health estimates}.

\end{CSLReferences}






\end{document}
